\chapter{Réduction des endomorphismes}
\origpage{339--358}

L'utilisation des démarches itératives, rendues encore plus importantes par l'informatique, conduit, dans le cadre des endomorphismes $u$ sur les espaces vectoriels $E$, à évaluer les puissances $u^n$ des opérateurs linéaires, ainsi que des puissances du type $(u+v)^n$. Mais dans ce cas, si $u$ et $v$ ne commutent pas, le calcul devient inextricable. Aussi est-il intéressant de savoir si on peut trouver des sous-espaces $F$ de $E$ sur lesquels $u$ commuterait avec tout opérateur linéaire de $F$, c'est-à-dire sur lequel $u$ agirait comme une homothétie, (nous allons justifier que, dans l'anneau $L(E,E)$, le centre est formé des homothéties). C'est cette démarche qui est à l'origine de l'introduction des notions de valeurs propres et de vecteurs propres.

Avant autre chose, justifions donc le

\medskip\noindent\textsc{Théorème 10.1.} --- \emph{Soit un espace vectoriel $E$ sur un corps $K$. Le centre de l'anneau $\operatorname{Hom}(E)$ des endomorphismes de $E$ est l'ensemble des homothéties.}

\subsection*{10.2.}
On appelle homothétie de rapport $\lambda\in K$, l'application linéaire $h_\lambda$ définie par :
\[
\forall x\in E,\qquad h_\lambda(x)=\lambda x.
\]
Soit $B=(e_i)_{i\in I}$ une base de $E$, (Théorème 6.71 pour l'existence).

Soit $f$ dans le centre de $\operatorname{Hom}(E)$, (c'est-à-dire l'ensemble des éléments qui commutent pour le produit de composition, avec tous les autres endomorphismes).

Pour $j\in I$ fixé, on décompose $f(e_j)$ en
\[
f(e_j)=\sum_{i\in I}\alpha_{ij}e_i,
\]
les $(\alpha_{ij})_{i\in I}$ étant presque tous nuls, si $\operatorname{card}I$ est infini, (Théorème 6.75).

Soit alors $g\in\operatorname{Hom}(E)$ définie par $g(e_j)=e_j$ et $\forall i\ne j$, $g(e_i)=0$, (on a vu qu'une application linéaire est déterminée par la donnée des images des vecteurs d'une base, théorème 6.77). Comme $f\circ g=g\circ f$, en particulier
\[
f\circ g(e_j)=f(e_j)=\sum_{i\in I}\alpha_{ij}e_i
=g(f(e_j))=\sum_{i\in I}\alpha_{ij}g(e_i)=\alpha_{jj}e_j
\]
d'où, (unicité de la décomposition d'un vecteur dans une base), $\forall i\ne j$, $\alpha_{ij}=0$, et donc $f(e_j)=\alpha_{jj}e_j$.

Si $\operatorname{card}I\ge2$, avec $j\ne k$ dans $I$ et $\alpha_{jj},\alpha_{kk}$ tels que $f(e_k)=\alpha_{kk}e_k$, en considérant $g$ définie par $g(e_j)=e_k$, $g(e_k)=e_j$ et $g(e_i)$ est n'importe quoi, 0 par exemple, pour $i$ dans $I-\{j,k\}$, on aura $(f\circ g)(e_j)=f(e_k)=\alpha_{kk}e_k$ et $(g\circ f)(e_j)=g(\alpha_{jj}e_j)=\alpha_{jj}e_k$, d'où $\alpha_{jj}=\alpha_{kk}$.

En notant $\lambda$ cette valeur commune de tous les scalaires $\alpha_{jj}$, on a donc, $\forall i\in I$, $f(e_i)=\lambda e_i$, résultat vrai si $\operatorname{card}I=1$.

Il est alors facile de voir que, $\forall x\in E$, $f(x)=\lambda x$, donc $f$ est bien une homothétie de rapport $\lambda$. Comme réciproquement, une homothétie $h_\lambda$ commute avec toute application linéaire, (vérification immédiate), on a bien le centre de $\operatorname{Hom}(E)$ constitué des homothéties. \bookend

\medskip\noindent\textsc{Remarque 10.3.} --- Il est parfois utile de savoir que $h$ linéaire de $E$ dans $E$ est une homothétie si et seulement si, $\forall x\in E$, $\exists\lambda_x\in K$, $h(x)=\lambda_x\cdot x$.

La justification du fait que $\lambda_x$ est indépendant de $x$ s'inspirant de la démonstration précédente.

\BellaSection{1. Valeurs spectrales, propres, vecteurs propres}

\medskip\noindent\textsc{Définition 10.4.} --- \emph{Soit $E$ un espace vectoriel sur un corps $K$ et $u$ un endomorphisme de $E$, appelé encore opérateur. Un scalaire $\lambda$ de $K$ est une valeur spectrale de $u$ si et seulement si \BellaCorrectionNote{$u-\lambda\,\mathrm{id}_E$}{$u=-\lambda\,\mathrm{id}_E$}{La condition porte sur la non-bijectivité de $u-\lambda\,\mathrm{id}_E$; l'égalité imprimée rendrait la définition incompatible avec la phrase suivante et avec tout le développement du chapitre.} est non bijective, et $\lambda\in K$ est une valeur propre de $u$ si et seulement si $u-\lambda\,\mathrm{id}_E$ est non injective.}

Toute valeur propre est donc spectrale, la réciproque étant vraie en dimension finie, (corollaire 6.85) car alors pour un endomorphisme les propriétés d'injectivité, de surjectivité et de bijectivité sont équivalentes.

\medskip\noindent\textsc{Définition 10.5.} --- \emph{Soit $u$ un opérateur linéaire \BellaCorrectionNote{sur un espace vectoriel sur $K$}{vectoriel sur $K$}{L'expression imprimée omet l'espace sur lequel agit l'opérateur; le contexte et la définition précédente imposent ce complément.}, on appelle spectre de $u$, l'ensemble $\operatorname{Sp}(u)$ des valeurs spectrales.}

Nous allons, dans ce chapitre, nous consacrer à l'étude des valeurs propres, celle des valeurs spectrales étant d'un domaine plus élevé.

\medskip\noindent\textsc{Définition 10.6.} --- \emph{Soit $u$ un opérateur linéaire sur $E$ vectoriel sur $K$, et $\lambda$ une valeur propre de $u$. On appelle sous-espace propre de $u$ pour la valeur propre $\lambda$ le noyau \BellaCorrectionNote{$\operatorname{Ker}(u-\lambda\,\mathrm{id}_E)$}{$\operatorname{Ker}(u-\mathrm{id}_E)$}{Le facteur $\lambda$ manque dans l'impression; sans lui le sous-espace défini ne dépendrait pas de la valeur propre annoncée.}. Les vecteurs de ce noyau sont les vecteurs propres associés à $\lambda$.}

On a donc :

\medskip\noindent\textsc{Remarque 10.7.} --- $(\lambda\text{ propre de }u)\Longleftrightarrow(\exists x\ne0,\ x\in E,\ u(x)=\lambda x)$.

Cette remarque est fondamentale. Par exemple, elle permet d'obtenir le résultat suivant : si
\[
P(X)=\sum_{n=0}^p a_nX^n
\]
est un polynôme de $K[X]$, tel que l'endomorphisme
\[
P(u)=\sum_{n=0}^p a_nu^n
\]
est nul, (avec $u^n=u\circ u\circ\cdots\circ u$, $n$ fois), les valeurs propres $\lambda$ de $u$ dans $K$ doivent être des racines de $P$, (pas forcément toutes les..., il y a une différence entre un article indéfini et un article défini) car, avec $x\ne0$ tel que $u(x)=\lambda x$ on vérifie aisément que $u^n(x)=\lambda^nx$, d'où
\[
P(u)(x)=0=\sum_{n=0}^pa_n\lambda^nx
\]
soit $P(\lambda)\cdot x=0$ et comme $x\ne0$, c'est que $P(\lambda)=0$.

\medskip\noindent\textsc{Théorème 10.8.} --- \emph{Soit $u$ un endomorphisme de $E$ vectoriel sur $K$. Des sous-espaces propres, $E_1,\ldots,E_n$ en nombre fini, associés à des valeurs propres distinctes, sont en somme directe.}

On va procéder par récurrence sur $n$, en commençant par vérifier la propriété pour $n=2$.

Soient $E_1$ et $E_2$ les noyaux de \BellaCorrectionNote{$u-\lambda_1\,\mathrm{id}$ et $u-\lambda_2\,\mathrm{id}$}{$u-\lambda\,\mathrm{id}$ et $u-\lambda\,\mathrm{id}$}{La phrase suivante utilise explicitement les deux valeurs propres distinctes $\lambda_1$ et $\lambda_2$; les indices ont été omis à l'impression.}. Il faut vérifier que $E_1\cap E_2=\{0\}$, (définition 6.45). Or si $x\in E_1\cap E_2$ on a $u(x)=\lambda_1x=\lambda_2x$, d'où $(\lambda_1-\lambda_2)x=0$ avec $\lambda_1\ne\lambda_2$ donc $x=0$ : on a bien $E_1$ et $E_2$ en somme directe.

On suppose le résultat vrai pour $n-1$ sous-espaces propres associés à des valeurs propres distinctes de $u$, et soient $\lambda_1,\ldots,\lambda_n$, $n$ valeurs propres distinctes de $u$, et $E_1,\ldots,E_n$ les sous-espaces propres associés. Pour justifier qu'ils sont en somme directe, il faut, (Théorème 6.51) vérifier que, $\forall i\in\{1,\ldots,n\}$,
\[
E_i\cap\left(\sum_{\substack{j=1\\j\ne i}}^nE_j\right)=\{0\}.
\]
Soit $x$ dans cette intersection, $x\in E_i$ d'où $u(x)=\lambda_ix$ et d'autre part
\[
x\in\sum_{\substack{j=1\\j\ne i}}^nE_j
\]
donc il existe des $x_j$ dans les $E_j$ tels que
\[
x=\sum_{\substack{j=1\\j\ne i}}^nx_j.
\]
Par linéarité,
\[
u(x)=\sum_{\substack{j=1\\j\ne i}}^nu(x_j)
\]
avec $u(x_j)=\lambda_jx_j$.

On a donc
\[
\lambda_ix=\sum_{\substack{j=1\\j\ne i}}^n\lambda_jx_j
\]
mais
\[
\lambda_ix=\lambda_i\sum_{\substack{j=1\\j\ne i}}^nx_j
\]
d'où l'égalité
\[
\sum_{\substack{j=1\\j\ne i}}^n(\lambda_i-\lambda_j)x_j=0.
\]
Mais comme les $(E_j)_{j=1\ldots n,\ j\ne i}$ sont en somme directe, (hypothèse de récurrence), chaque $(\lambda_i-\lambda_j)x_j$ est nul (Théorème 6.51) et comme $\lambda_i-\lambda_j\ne0$, (valeurs propres distinctes) c'est que chaque $x_j$ est nul d'où finalement $x=0$. \bookend

Le plus souvent, cette somme directe ne donne pas l'espace entier d'où la définition suivante :

\medskip\noindent\textsc{Définition 10.9.} --- \emph{Un endomorphisme $u$ sur un espace vectoriel $E$ est dit diagonalisable si et seulement si $E$ est somme directe d'un nombre fini de \BellaCorrection{sous-espaces propres}{sous-espace propres}.}

Comme les sous-espaces propres sont forcément en somme directe, cette définition ne fait que préciser le fait que cette somme directe est l'espace entier.

Elle n'est pas pleinement satisfaisante. En effet, on peut remarquer que si $u$ est diagonalisable, avec
\[
E=\bigoplus_{i=1}^nE_i,
\]
où $E_i=\operatorname{Ker}(u-\lambda_i\mathrm{id})$ est le sous-espace propre associé à la valeur propre $\lambda_i$, si $B_i$ est une base de $E_i$,
\[
B=\bigcup_{i=1}^nB_i
\]
est une base de $E$, (Théorème 6.81) formée de vecteurs propres, et pour beaucoup d'applications on aurait pu prendre comme définition :
\[
(u\text{ diagonalisable})\Longleftrightarrow(\text{existence d'une base de vecteurs propres}).
\]
Mais ceci ne conviendrait pas pour les polynômes d'endomorphismes... s'il y a une infinité de valeurs propres distinctes.

Avant d'abandonner ce sujet, un dernier mot : ce cas existe : prendre $E$ vectoriel dénombrable, $B=(e_n)_{n\in\mathbb N}$ une base et définir l'endomorphisme $u$ par $u(e_n)=ne_n$ par exemple : la base $B$ est une base de vecteurs propres et comme on a une infinité de valeurs propres distinctes, on ne peut pas avoir $E$ somme directe d'un nombre fini de sous-espaces propres distincts. Pour réaliser la situation décrite, on peut prendre $E=\mathbb R[X]$ et définir $u$ par : $P(X)\longmapsto Q(X)$ avec $Q(X)=XP'(X)$.

Travaillons dans ce qui suit avec la définition 10.9. Le lecteur curieux pourra vérifier quels résultats restent valables avec la 2ème définition.

\medskip\noindent\textsc{Théorème 10.10.} --- \emph{Soit $u$ un endomorphisme diagonalisable de $E$ et $F$ un sous-espace de $E$, stable par $u$. Alors l'endomorphisme $\bar u$ induit par $u$ sur $F$ est aussi diagonalisable.}

\subsection*{10.11.}
Rappelons que $\bar u$ est défini de $F$ dans $F$ par $\bar u(x)=u(x)$ alors que la restriction, $u|_F$ serait l'application de $F$ dans $E$ définie par $u|_F(x)=u(x)$.

Soient $E_1,\ldots,E_n$ les sous-espaces propres distincts de $u$, associés aux valeurs propres distinctes $\lambda_1,\ldots,\lambda_n$ de $u$, tels que
\[
E=\bigoplus_{i=1}^nE_i.
\]
Soit $x\in F$, comme $x\in E$, $\exists$ des $x_i$ dans les $E_i$, tels que
\[
x=\sum_{i=1}^nx_i.
\]
En fait les $x_i$ sont aussi dans $F$. En prenant les images successives de l'égalité $x=x_1+x_2+\cdots+x_n$ par $u$, et comme $u(x_i)=\lambda_ix_i$ on obtient le système
\[
\tag{I}
\left\{
\begin{aligned}
x&=x_1+x_2+\cdots+x_n,\\
u(x)&=\lambda_1x_1+\lambda_2x_2+\cdots+\lambda_nx_n,\\
u^2(x)&=(\lambda_1)^2x_1+\cdots+(\lambda_n)^2x_n,\\
&\ \vdots\\
u^{n-1}(x)&=(\lambda_1)^{n-1}x_1+\cdots+(\lambda_n)^{n-1}x_n,
\end{aligned}
\right.
\]
qui se \BellaCorrection{résout}{résoud} en :
\[
\tag{II}
x_i=
\frac{
\begin{vmatrix}
1&\cdots&1&x&1&\cdots&1\\
\lambda_1&\cdots&\lambda_{i-1}&u(x)&\lambda_{i+1}&\cdots&\lambda_n\\
\vdots&&\vdots&\vdots&\vdots&&\vdots\\
\lambda_1^{n-1}&\cdots&\lambda_{i-1}^{n-1}&u^{n-1}(x)&\lambda_{i+1}^{n-1}&\cdots&\lambda_n^{n-1}
\end{vmatrix}
}{\det V(\lambda_1,\ldots,\lambda_n)},
\]
le numérateur étant un « déterminant vecteur » que l'on développe suivant la $i$ème colonne. (Voir Ch. 9, exercice 6 pour Van der Monde).

(Pour justifier ceci, si $B=(e_k)_{k\in K}$ est une base de $E$, et si les $\mu_k$ sont les formes coordonnées associées dans le dual, (voir 6.109) en prenant l'image de (I) par $\mu_k$ on obtient un système linéaire classique qui donnera la $k$ième coordonnée $\mu_k(x_i)$ par la valeur obtenue dans (II) en remplaçant la colonne de vecteurs par leurs $k$ième coordonnées respectives, voir 9.55).

Il en résulte que
\[
x_i\in\operatorname{Vect}(x,u(x),u^2(x),\ldots,u^{n-1}(x))\subset F
\]
car $F$ est stable par $u$ et $x\in F$. Donc
\[
x=\sum_{i=1}^nx_i\qquad\text{avec }x_i\in F\cap E_i,
\]
d'où
\[
F=\sum_{i=1}^n(F\cap E_i).
\]
Cette somme est directe car
\[
(F\cap E_i)\cap\left(\sum_{\substack{j=1\\j\ne i}}^n(F\cap E_j)\right)
\subset E_i\cap\left(\sum_{\substack{j=1\\j\ne i}}^nE_j\right)=\{0\}.
\]
D'où $\bar u$ est diagonalisable. \bookend

\medskip\noindent\textsc{Remarque 10.12.} --- Quand on écrit
\[
F=\bigoplus_{j=1}^n(F\cap E_j),
\]
certains $F\cap E_j$ peuvent être réduits à $\{0\}$.

\medskip\noindent\textsc{Remarque 10.13.} --- Le résultat précédent, (et le suivant) restent valables avec la définition élargie de diagonalisable.

\medskip\noindent\textsc{Théorème 10.14.} --- \emph{Soient $u$ et $v$ deux endomorphismes qui commutent. Tout sous-espace propre de l'un est stable par l'autre.}

En effet, si $\lambda$ est valeur propre de $u$, et si $x\in\operatorname{Ker}(u-\lambda\mathrm{id})$ on a
\[
u(v(x))=v(u(x))=v(\lambda x)=\lambda v(x)
\]
car $v$ est linéaire donc $(u-\lambda\mathrm{id})(v(x))=0$, soit $v(x)\in\operatorname{Ker}(u-\lambda\mathrm{id})$. \bookend

\medskip\noindent\textsc{Corollaire 10.15.} --- \emph{Si $u$ et $v$ sont deux endomorphismes de $E$ diagonalisables qui commutent, il existe une base commune de diagonalisation.}

Car si
\[
E=\bigoplus_{i=1}^nE_i
\]
avec les $E_i$ sous-espaces propres de $u$, comme $u$ et $v$ commutent, $E_i$ est stable par $v$, et comme $v$ est diagonalisable, $v_i$ induit par $v$ sur $E_i$ est diagonalisable, (10.10) donc il existe une base $B_i$ de $E_i$, formée de vecteurs propres pour $v_i$, donc pour $v$. Comme on est dans $E_i$, sous-espace propre de $u$, tous les vecteurs de $B_i$ sont aussi vecteurs propres pour $u$, et finalement
\[
B=\bigcup_{i=1}^nB_i
\]
est une base de vecteurs propres pour $u$, et pour $v$. \bookend

On appelle \emph{base de diagonalisation} pour $u$ diagonalisable, toute base formée de vecteurs propres.

Ces premiers résultats sur les vecteurs propres montrent l'importance de la commutativité entre opérateurs linéaires, et de la notion de stabilité des sous-espaces. Avant de traiter le cas des espaces vectoriels de dimension finie, justifions le

\medskip\noindent\textsc{Théorème 10.16.} --- \emph{Un hyperplan $H$ de $E$ vectoriel est stable par un endomorphisme $u$ si et seulement si $H$ est le noyau d'une forme linéaire, $\varphi$, vecteur propre non nul de ${}^tu$.}

Rappelons que si $u\in L(E,F)$, ${}^tu\in L(F^*,E^*)$ est défini par ${}^tu(\varphi)=\varphi\circ u$, si $\varphi$ est dans le dual $F^*$ de $F$, (définition 6.100).

Soit $H$ un hyperplan de $E$, stable par $u$ et $\varphi$ dans $E^*$ une forme linéaire ayant $H$ pour noyau. (Si $\{a\}$ est une base d'un supplémentaire de $H$ dans $E$, et $C$ une base de $H$, $B=C\cup\{a\}$ est une base de $E$ et on peut définir $\varphi$ par $\varphi(a)=1$ et $\varphi$ (tout vecteur de $C$) $=0$).

Alors, $\forall x\in H$, $u(x)\in H$ donc $\varphi(x)=0\Rightarrow\varphi(u(x))=0$. Avec $E=H\oplus K a$, ($K$ corps de base) on a $\varphi(a)\ne0$, sinon $\operatorname{Ker}\varphi=E$, donc
\[
{}^tu(\varphi)(a)=\varphi(u(a))=\frac{\varphi(u(a))}{\varphi(a)}\varphi(a).
\]
Les deux formes linéaires ${}^tu(\varphi)$ et
\[
\frac{\varphi(u(a))}{\varphi(a)}\varphi
\]
prennent donc la même valeur sur $a$, mais aussi sur $H$, et là, elles s'annulent, donc elles sont égales : avec
\[
\lambda=\frac{\varphi(u(a))}{\varphi(a)}
\]
on a ${}^tu(\varphi)=\lambda\varphi$ : $\varphi$ est vecteur propre pour la valeur propre $\lambda$ de ${}^tu$.

\emph{Réciproquement} soit $\varphi$ un vecteur propre, dans $E^*$, de ${}^tu$, et $H=\operatorname{Ker}\varphi$, avec $\varphi\ne0$ bien sûr, alors l'hyperplan $H$ est stable pour $u$ car si $x\in H$, on a :
\[
\varphi(u(x))={}^tu(\varphi)(x)=(\lambda\varphi)(x)=\lambda\cdot\varphi(x)=0
\]
car $x\in\operatorname{Ker}\varphi$.

Ici $\lambda$ est la valeur propre associée à $\varphi$ vecteur propre de ${}^tu$. \bookend

\BellaSection{2. Étude en dimension finie}

Soit $E$ espace vectoriel de dimension finie, $n$, sur un corps $K$. On sait alors que, pour un endomorphisme de $E$, les notions d'injectivité, de surjectivité ou de bijectivité sont équivalentes (corollaire 6.85).

Donc le scalaire $\lambda$ de $K$ sera valeur propre de $u$ si et seulement si c'est une valeur spectrale, (définition 10.4) et c'est pourquoi, en dimension finie, le spectre d'un opérateur devient l'ensemble de ses valeurs propres.

Soit alors $B=\{e_1,\ldots,e_n\}$ une base de $E$ et $U$ la matrice de $u$ dans cette base. Celle de $u-\lambda\mathrm{id}_E$ est $U-\lambda I_n$, $I_n$ matrice diagonale unité d'ordre $n$, et
\[
(\lambda\text{ valeur propre de }u)\Longleftrightarrow(\det(U-\lambda I_n)=0),
\]
puisque les matrices non inversibles, donc associées à des applications non bijectives, sont celles de déterminant nul, voir Théorème 9.42.

Cette expression $\det(U-\lambda I_n)$ est indépendante du choix de la base $B$ car si $B'$ est une autre base, et si $P$ est la matrice de passage de $B$ à $B'$, la matrice $U'$ de $u$ dans la nouvelle base $B'$ est $U'=P^{-1}UP$, (8.23 et 8.25). Elle conduit à la relation :
\[
(\lambda\text{ valeur propre de }u)\Longleftrightarrow(\det(U'-\lambda I_n)=0),
\]
or
\[
U'-\lambda I_n=P^{-1}UP-\lambda P^{-1}I_nP=P^{-1}(U-\lambda I_n)P
\]
d'où
\[
\det(U'-\lambda I_n)=\det(P^{-1}(U-\lambda I_n)P)=\det P^{-1}\det(U-\lambda I_n)\det P
\]
et comme $\det P^{-1}\det P=1$ dans $K$, commutatif, il reste
\[
\det(U'-\lambda I_n)=\det(U-\lambda I_n).
\]

Cette expression est un polynôme en $\lambda$, de degré $n$.

En effet, si $u_{ij}$ désigne le terme général, ($i$ème ligne, $j$ème colonne) de $U$, en notant $V=U-\lambda I_n$ on aura $v_{ii}=u_{ii}-\lambda$ et si $i\ne j$, $v_{ij}=u_{ij}$ donc $v_{ij}$ est un polynôme de degré 1 ou 0 en $\lambda$, et en utilisant l'expression développée d'un déterminant, on constate que
\[
\det(U-\lambda I_n)=\det V=\sum_{\sigma\in\mathfrak S_n}\varepsilon_\sigma
v_{1,\sigma(1)}v_{2,\sigma(2)}\cdots v_{n,\sigma(n)}
\]
est une somme de $n!$ termes, (indexés par les éléments du groupe symétrique des permutations de $\{1,\ldots,n\}$), chaque terme étant un polynôme en $\lambda$ de degré $n$ au plus : on a bien un polynôme de degré $n$ au plus.

On peut facilement préciser les coefficients des termes de degré $n$ et $n-1$. En effet, si $\sigma\ne\mathrm{id}$, $\exists i\in\{1,\ldots,n\}$ tel que $\sigma(i)=j$ avec $j\ne i$. Mais alors dans le terme $\varepsilon_\sigma v_{1\sigma(1)}\cdots v_{i\sigma(i)}\cdots v_{n\sigma(n)}$ le facteur $v_{ij}=u_{ij}$ est de degré 0 en $\lambda$, le produit est de degré $n-1$ au plus; mais en fait $v_{ij}$ est aussi la contribution de la $j$ème colonne : comme $\sigma$ est une bijection on ne peut pas avoir $\sigma(j)=j$, donc le terme $v_{j\sigma(j)}$ ne peut être $v_{jj}$, il est du type $v_{jk}$ avec $k\ne j$, donc lui aussi de degré 0 en $\lambda$.

Finalement, si $\sigma\ne\mathrm{id}$, $\varepsilon_\sigma v_{1\sigma(1)}\cdots v_{n\sigma(n)}$ est de degré $n-2$ au plus. Les termes de degré $n$ et $n-1$ du polynôme $\det(U-\lambda I_n)$ proviennent donc de
\[
\varepsilon_{\mathrm{id}}v_{11}v_{22}\cdots v_{nn}=(u_{11}-\lambda)(u_{22}-\lambda)\cdots(u_{nn}-\lambda)
\]
et valent donc $(-\lambda)^n$ et $(-\lambda)^{n-1}(u_{11}+u_{22}+\cdots+u_{nn})$ respectivement.

\subsection*{10.17.}
On sait que
\[
\sum_{i=1}^nu_{ii}
\]
est ce qu'on appelle la trace de la matrice $U$, ou encore de l'endomorphisme $u$, puisque le polynôme étant indépendant de la base utilisée, on aura le même coefficient de $\lambda^{n-1}$ dans une autre base.

Enfin, pour $\lambda=0$, $\det(U-\lambda I_n)=\det U$, et c'est le terme constant du polynôme. On a finalement obtenu le résultat suivant :

\medskip\noindent\textsc{Théorème 10.18.} --- \emph{Soit $u$ un endomorphisme de $E\simeq K^n$. Les valeurs propres $\lambda$ de $u$ sont les racines, dans $K$, du polynôme noté $\chi_u(\lambda)$ défini par $\chi_u(\lambda)=\det(U-\lambda I_n)$, $U$ étant la matrice de $u$ dans une base $B$ de $E$, et ce polynôme étant indépendant de la base.}

\medskip\noindent\textsc{Définition 10.19.} --- \emph{Ce polynôme $\chi_u(\lambda)=\det(U-\lambda I_n)$ est le polynôme caractéristique de l'endomorphisme $u$.}

C'est une notion attachée aux espaces vectoriels de dimension finie.

On a
\[
\chi_u(\lambda)=(-1)^n\lambda^n+(-1)^{n-1}\operatorname{trace}(u)\lambda^{n-1}+\cdots+\det u.
\]
Certains préfèrent prendre pour polynôme caractéristique le polynôme unitaire $\det(\lambda I_n-U)$, ce n'est qu'une question de convention.

Du fait de l'existence de ce polynôme caractéristique, il faut comprendre qu'une grande injustice se prépare. En dimension infinie, le fait d'avoir des valeurs propres va être une faveur accordée à certains endomorphismes (allez donc trouver une valeur propre de $u:\mathbb R[X]\longmapsto\mathbb R[X]$ défini par $P\longmapsto Q$ avec $Q(X)=XP(X)$), alors qu'en dimension finie pour peu \BellaCorrection{que le corps}{que que le corps} soit algébriquement clos, ($\mathbb C$ par exemple), tous les endomorphismes auront droit à leurs valeurs propres, racines de leur polynôme caractéristique !

Voyons maintenant comment chercher les vecteurs propres.

On suppose toujours une base $B=\{e_1,\ldots,e_n\}$ de $E$ fixée, et $U$ désigne la matrice de $u$ dans la base $B$.

Soit $\lambda$ une racine du polynôme caractéristique $\chi_u(\lambda)=\det(U-\lambda I_n)$. La matrice $U-\lambda I_n$ est non injective, et son noyau est le sous-espace propre de $u$ pour la valeur propre $\lambda$.

Il est donc de dimension $n-\operatorname{rang}(U-\lambda I_n)$ (corollaire 6.85).

Donc, connaissant $\lambda$, on déterminera $r=\operatorname{rang}(U-\lambda I_n)$, pour connaître la dimension du sous-espace propre $\operatorname{Ker}(u-\lambda\mathrm{id}_E)$, et si on veut que
\[
x=\sum_{i=1}^nx_ie_i
\]
soit vecteur propre, on doit résoudre le système linéaire
\[
(U-\lambda I_n)
\begin{pmatrix}x_1\\x_2\\\vdots\\x_n\end{pmatrix}=0
\]
pour connaître la matrice $X$ des composantes de $x$ dans la base $B$.

Voir en 9.47 et suivants, la résolution d'un tel système.

\medskip\noindent\textbf{Quelques propriétés du polynôme caractéristique et des valeurs propres}

\medskip\noindent\textsc{Remarque 10.20.} --- On peut calculer le coefficient de $\lambda^{n-2}$. En effet, en reprenant le calcul du début du paragraphe pour $\sigma\ne\mathrm{id}$, dans
\[
\varepsilon_\sigma v_{1\sigma(1)}\cdots v_{n\sigma(n)},
\]
avec $i$ tel que $j=\sigma(i)$ soit $\ne i$ on a vu que $v_{ij}=u_{ij}$ est de degré 0 en $\lambda$ ainsi que $v_{j\sigma(j)}=v_{jk}$ avec $k\ne j$.

Mais si $\sigma(j)\ne i$, (et $\ne j$ aussi, sinon, $\sigma(j)=j\Rightarrow j=\sigma(j)=\sigma(i)$ exclu car $\sigma$ bijection) en posant $\sigma(j)=k$, on a le terme $u_{jk}$ qui est la contribution de la $k$ème colonne au produit, donc $u_{kk}-\lambda$ ne figurera pas dans le produit qui sera de degré $n-3$ au plus.

Les termes en $\lambda^{n-2}$ proviennent donc des transpositions et de l'identité. Pour une transposition $T_{ij}$ on a
\[
-u_{ij}u_{ji}\prod_{\substack{k\ne i\\k\ne j}}(u_{kk}-\lambda)
\]
qui donne pour terme en $\lambda^{n-2}$ :
\[
-(-\lambda)^{n-2}u_{ij}u_{ji}.
\]
Pour $\sigma=\mathrm{id}$, dans $(u_{11}-\lambda)\cdots(u_{nn}-\lambda)$ le terme en $\lambda^{n-2}$ s'obtient en prenant le terme constant dans 2 facteurs : c'est donc
\[
\left(\sum_{1\le i<j\le n}u_{ii}u_{jj}\right)(-\lambda)^{n-2},
\]
d'où
\[
\lambda^{n-2}(-1)^{n-2}
\left(\sum_{1\le i<j\le n}(u_{ii}u_{jj}-u_{ij}u_{ji})\right)
\]
comme terme en $\lambda^{n-2}$.

On pourrait poursuivre ainsi en introduisant les cycles de 3 termes pour chercher le terme en $\lambda^{n-3}$.

\medskip\noindent\textsc{Théorème 10.21.} --- \emph{Soient deux endomorphismes $u$ et $v$ de $E\simeq K^n$ on a $\chi_{uv}(\lambda)=\chi_{vu}(\lambda)$.}

Il y a plusieurs procédés de justification. En voici un. On fixe une base $B$ de $E$ d'où des matrices $U$ et $V$ associées à $u$ et $v$.

En effectuant des produits matriciels par blocs, on constate que
\[
\begin{pmatrix}
VU-\lambda I_n&-V\\
0&-\lambda I_n
\end{pmatrix}
\begin{pmatrix}
I_n&0\\
U&I_n
\end{pmatrix}
=
\begin{pmatrix}
-\lambda I_n&-V\\
-\lambda U&-\lambda I_n
\end{pmatrix}
\]
alors que
\[
\begin{pmatrix}
I_n&0\\
U&I_n
\end{pmatrix}
\begin{pmatrix}
-\lambda I_n&-V\\
0&UV-\lambda I_n
\end{pmatrix}
=
\begin{pmatrix}
-\lambda I_n&-V\\
-\lambda U&-\lambda I_n
\end{pmatrix},
\]
d'où en égalant les déterminants de ces deux produits, l'égalité
\[
(-\lambda)^n\det(VU-\lambda I_n)=(-\lambda)^n\det(UV-\lambda I_n)
\]
donc
\[
\det(VU-\lambda I_n)=\det(UV-\lambda I_n)
\]
soit $\chi_{vu}(\lambda)=\chi_{uv}(\lambda)$. \bookend

\medskip\noindent\textsc{Théorème 10.22.} --- \emph{Dit d'Hadamard, pour la localisation des valeurs propres. --- Soit une matrice $U$ de terme général $u_{ij}$ complexe. Les valeurs propres de $U$ sont dans la réunion des disques $D_i$ du plan complexe, centrés en $u_{ii}$, de rayon}
\[
\Lambda_i=\sum_{\substack{j=1\\j\ne i}}^n|u_{ij}|.
\]

En effet soit $\lambda$ valeur propre de $u$ et $Z$ un vecteur propre associé, assimilé à une matrice colonne. On choisit bien sûr $Z\ne0$, donc $\exists i$ tel que $|z_i|\ne0$. En prenant les $i$èmes composantes des 2 termes de l'égalité $UZ=\lambda Z$ il vient
\[
\sum_{j=1}^nu_{ij}z_j=\lambda z_i,
\]
soit encore
\[
(\lambda-u_{ii})z_i=\sum_{\substack{j=1\\j\ne i}}^nu_{ij}z_j,
\]
d'où en divisant par $z_i\ne0$ et en utilisant l'inégalité triangulaire :
\[
|\lambda-u_{ii}|\le\sum_{\substack{j=1\\j\ne i}}^n|u_{ij}|\left|\frac{z_j}{z_i}\right|.
\]
Si on a pris la précaution de choisir $i$ tel que, non seulement $z_i\ne0$, mais en fait tel que $|z_i|$ soit le sup de $\{|z_k|;\ k=1,\ldots,n\}$, chaque $|z_j/z_i|$ se majore par 1 et il vient
\[
|\lambda-u_{ii}|\le\sum_{\substack{j=1\\j\ne i}}^n|u_{ij}|=\Lambda_i,
\]
donc $\lambda$ est dans le disque $D_i$ de centre $u_{ii}$, de rayon $\Lambda_i$ et on a bien trouvé un indice $i$ tel que $\lambda\in D_i$, d'où
\[
\operatorname{Sp}(U)\subset\bigcup_{i=1}^nD_i,
\]
$\operatorname{Sp}(U)$ désignant le spectre de $U$. \bookend

\medskip\noindent\textsc{Corollaire 10.23.} --- \emph{On a aussi $\operatorname{Sp}(U)\subset\bigcup_{j=1}^n\Delta_j$, avec $\Delta_j$ disque complexe de centre $u_{jj}$ et de rayon}
\[
\mu_j=\sum_{\substack{i=1\\i\ne j}}^n|u_{ij}|.
\]

Comme $\det({}^tA)=\det A$, le spectre de la matrice ${}^tU$ est celui de la matrice $U$. Mais en appliquant Hadamard à ${}^tU$, les disques $D_j$ pour ${}^tU$ sont les $\Delta_j$. \bookend

\medskip\noindent\textsc{Corollaire 10.24.} --- \emph{Si une matrice carrée $U$ complexe est telle que}
\[
\forall i=1,\ldots,n,\qquad |u_{ii}|>\sum_{\substack{j=1\\j\ne i}}^n|u_{ij}|,
\]
\emph{elle est inversible.}

Car chaque $D_i$ ne contient pas 0, donc $0\notin\bigcup_iD_i$ : 0 ne peut pas être valeur propre, d'où $U$ injective, donc inversible. \bookend

\subsection*{10.25.}
Une telle matrice, où, pour chaque ligne, le module du terme diagonal est strictement supérieur à la somme des modules des coefficients non diagonaux de la ligne est dite à \emph{diagonale fortement dominante}.

On procède de même pour les colonnes.

Nous allons maintenant essayer de réduire un endomorphisme à une forme plus simple, dans le cas de la dimension finie. On a déjà :

\medskip\noindent\textsc{Théorème 10.26.} --- \emph{Soit $u$ un endomorphisme de $E\simeq K^n$, de polynôme caractéristique $\chi_u(X)$. Si $\lambda$ est une racine de $\chi_u$ de multiplicité $\alpha$, la dimension du sous-espace propre \BellaCorrection{associé}{associée} est inférieure ou égale à $\alpha$.}

Comme $\chi_u(\lambda)=0$, $E_\lambda=\operatorname{Ker}(u-\lambda\mathrm{id}_E)$ est non réduit à 0. Soit $C$ une base de $E_\lambda$ que l'on complète en $B=C\cup D$ base de $E$. On pose $p=\dim E_\lambda$. La matrice de $u$ dans la base $B$ est du type :
\[
U=
\begin{pmatrix}
\lambda I_p&A\\
0&B_0
\end{pmatrix},
\]
les $p$ premiers vecteurs étant ceux de $C$, les $n-p$ derniers ceux de $D$.

Donc, vu l'indépendance du choix de la base pour calculer le polynôme caractéristique,
\[
\chi_u(X)=\det(U-XI_n)
=\det\begin{pmatrix}
(\lambda-X)I_p&A\\
0&B_0-XI_{n-p}
\end{pmatrix}
=(\lambda-X)^p\det(B_0-XI_{n-p});
\]
il admet $\lambda$ comme racine avec une multiplicité $\ge p$. Comme on a appelé $\alpha$ cette multiplicité, on a donc $p\le\alpha$, soit encore
\[
\dim\operatorname{Ker}(u-\lambda\mathrm{id}_E)\le\text{multiplicité de la valeur propre}.
\]
\bookend

Il faut remarquer que l'inégalité peut être stricte. Par exemple la matrice
\[
U=
\begin{pmatrix}
a&1&0&\cdots&0\\
0&a&1&\ddots&\vdots\\
\vdots&\ddots&\ddots&\ddots&0\\
0&\cdots&0&a&1\\
0&\cdots&\cdots&0&a
\end{pmatrix}
\]
est celle d'un endomorphisme $u$ n'ayant que $a$ pour valeur propre, et
\[
\chi_u(\lambda)=
\begin{vmatrix}
a-\lambda&1&0&\cdots&0\\
0&a-\lambda&1&\ddots&\vdots\\
\vdots&\ddots&\ddots&\ddots&0\\
0&\cdots&0&a-\lambda&1\\
0&\cdots&\cdots&0&a-\lambda
\end{vmatrix}
=(a-\lambda)^n
\]
admet $a$ pour seul zéro, de multiplicité $n$.

Or
\[
U-aI_n=
\begin{pmatrix}
0&1&0&\cdots&0\\
0&0&1&\ddots&\vdots\\
\vdots&\ddots&\ddots&\ddots&0\\
0&\cdots&0&0&1\\
0&\cdots&\cdots&0&0
\end{pmatrix}
\]
est de rang $n-1$, le sous-espace propre est de dimension $n-(n-1)=1$, avec $1<n$.

Dans ce cas l'endomorphisme n'est pas diagonalisable. Plus précisément on a :

\medskip\noindent\textsc{Théorème 10.27.} --- \emph{Soit un endomorphisme $u$ sur $E\simeq K^n$. Il est diagonalisable si et seulement si son polynôme caractéristique a toutes ses valeurs propres dans le corps $K$ et si pour chaque valeur propre, la multiplicité est la dimension du sous-espace propre.}

Si $u$ est diagonalisable, avec $\lambda_1,\ldots,\lambda_p$ valeurs propres distinctes de $u$ et $E_1,\ldots,E_p$ les sous-espaces propres associés de dimensions respectives $r_1,\ldots,r_p$, en notant $B_i$ une base de $E_i$, dans la base $B=B_1\cup\cdots\cup B_p$ (vecteurs pris dans cet ordre), la matrice de $u$ est
\[
U=\operatorname{diag}(\underbrace{\lambda_1,\ldots,\lambda_1}_{r_1},
\underbrace{\lambda_2,\ldots,\lambda_2}_{r_2},\ldots,
\underbrace{\lambda_p,\ldots,\lambda_p}_{r_p})
\]
d'où
\[
\chi_u(\lambda)=\prod_{j=1}^p(\lambda_j-\lambda)^{r_j}:
\]
le polynôme caractéristique a toutes ses racines dans $K$ et, pour chacune d'entre elles, la multiplicité, $r_j$, est la dimension du sous-espace propre associé.

Réciproquement on suppose que $\chi_u(\lambda)$, de degré $n$, a toutes ses racines dans $K$ et, si $\lambda_1,\ldots,\lambda_p$ sont les racines distinctes, de multiplicités respectives $\alpha_1,\ldots,\alpha_p$, on suppose que pour chaque $j\in\{1,\ldots,p\}$,
\[
\dim(\operatorname{Ker}(u-\lambda_j\mathrm{id}))=\alpha_j.
\]
On sait déjà que les $E_j=\operatorname{Ker}(u-\lambda_j\mathrm{id})$ sont en somme directe, (Théorème 10.8), et
\[
F=\bigoplus_{j=1}^pE_j
\]
est de dimension
\[
\sum_{j=1}^p\alpha_j=\deg(\chi_u)=n,
\]
donc $F=E$ : on a bien $E$ somme directe des sous-espaces propres de $u$ qui est donc diagonalisable. \bookend

\medskip\noindent\textsc{Définition 10.28.} --- \emph{Un polynôme $P\in K[X]$ est dit scindé sur $K$ si et seulement si toutes ses racines sont dans $K$.}

Ce qui revient à dire qu'il est du type
\[
P(X)=k\prod_{i=1}^p(X-\lambda_i)^{\alpha_i}
\]
avec $k\in K^*$, les $\lambda_i\in K$, et $\alpha_1+\alpha_2+\cdots+\alpha_p=\deg(P)$. On prend $k\ne0$ car cette notion n'a pas d'intérêt pour le polynôme nul.

\medskip\noindent\textsc{Corollaire 10.29.} --- \emph{Soit $u$ un endomorphisme de $E\simeq K^n$, dont le polynôme caractéristique n'a que des racines simples toutes dans $K$ alors $u$ est diagonalisable.}

En effet, en notant $\lambda_1,\ldots,\lambda_n$ les $n$ valeurs propres distinctes, on a
\[
\chi_u(\lambda)=\prod_{i=1}^n(\lambda_i-\lambda)
\]
est scindé sur $K$ et, $\forall i$, $\dim\operatorname{Ker}(u-\lambda_i\mathrm{id}_E)\ge1$ (car $\lambda_i$ valeur propre) donc cette dimension est $\ge$ multiplicité de la racine. Comme elle est toujours $\le$ la multiplicité, (Théorème 10.26), elle lui est égale donc $u$ est diagonalisable. \bookend

Cette condition, suffisante, n'est pas nécessaire : une homothétie est diagonalisable bien que n'ayant qu'une valeur propre, de multiplicité $n$.

Et si on ne peut pas diagonaliser... , c'est fini? Pas forcément, on peut parfois trigonaliser.

\medskip\noindent\textsc{Définition 10.30.} --- \emph{Un endomorphisme $u$ de $E\simeq K^n$, est dit trigonalisable s'il existe une base $B$ dans laquelle la matrice $U$ de $u$ est triangulaire supérieure ou inférieure, peu importe.}

On parle de même d'une matrice trigonalisable $U$ si et seulement si il existe $P$ régulière telle que $P^{-1}UP$ soit triangulaire. Les 2 notions coïncident, si on prend $E=K^n$, rapporté à une base, par exemple la base canonique, et si $u$ est l'endomorphisme de matrice $U$ dans cette base.

Vérifions que trigonale inférieure ou supérieure, peu importe. Soit $B$ base dans laquelle la matrice de $u$ est $T$, triangulaire supérieure. En notant $t_{rs}$ le terme général de $T$, (donc $t_{rs}=0$ si $r>s$), et $\{e_1,\ldots,e_n\}$ la base $B$, on a, $\forall s$, $u(e_s)\in\operatorname{Vect}(e_1,e_2,\ldots,e_s)$.

Soit $C$ la base $B$ numérotée dans l'ordre $e_n,e_{n-1},\ldots,e_1$. Si on note $\varepsilon_k$ le vecteur générique de $C$, on a donc $\varepsilon_k=e_{n-k+1}$.

Mais alors
\[
u(\varepsilon_k)=u(e_{n-k+1})\in\operatorname{Vect}(e_1,e_2,\ldots,e_{n-k+1}),
\]
or
\[
e_1=\varepsilon_n,
\qquad e_2=\varepsilon_{n-1},\ldots,e_{n-k+1}=\varepsilon_k
\]
d'où $u(\varepsilon_k)\in\operatorname{Vect}(\varepsilon_k,\varepsilon_{k+1},\ldots,\varepsilon_n)$ : la matrice de $u$ dans la base $C$ est donc bien triangulaire inférieure. On procède de même si on part d'une matrice triangulaire inférieure. \bookend

\medskip\noindent\textsc{Théorème 10.31.} --- \emph{Soit $u$ un endomorphisme sur $E\simeq K^n$. Il est trigonalisable si et seulement si son polynôme caractéristique est scindé. Même énoncé pour une matrice $U$.}

Si $u$ est trigonalisable, il existe une base $B$ dans laquelle la matrice $U$ de $u$ est triangulaire (supérieure par exemple), donc du type
\[
U=
\begin{pmatrix}
u_{11}&u_{12}&u_{13}&\cdots&u_{1n}\\
0&u_{22}&u_{23}&\cdots&u_{2n}\\
\vdots&\ddots&\ddots&\ddots&\vdots\\
0&\cdots&0&u_{n-1,n-1}&u_{n-1,n}\\
0&\cdots&\cdots&0&u_{nn}
\end{pmatrix}
\]
d'où
\[
\chi_u(\lambda)=\prod_{i=1}^n(u_{ii}-\lambda):
\]
le polynôme caractéristique est scindé.

\emph{Réciproquement}, on procède par récurrence sur $n$. Si $n=1$, toute base de $E=K$ est telle que la \BellaCorrectionNote{matrice $(u_{11})$}{matrice $(1,1)$}{En dimension 1, la matrice d'un endomorphisme est une matrice $1\times1$ dont l'unique coefficient dépend de $u$; la notation imprimée $(1,1)$ ne peut être sa matrice en général.} de $u$ dans cette base est triangulaire.

On suppose le résultat vrai si $\dim E=n-1$, ou pour une matrice carrée d'ordre $n-1$.

Soit $u$ endomorphisme de $E\simeq K^n$, de matrice $U$ dans une base $B$ alors la matrice du transposé de $u$, ${}^tu$, endomorphisme de $E^*$, dans la base duale est ${}^tU$, (définition 8.8), donc $\chi_{{}^tu}=\chi_u$ est scindé sur $K$. (On a
\[
\chi_{{}^tu}(\lambda)=\det({}^tU-\lambda I_n)=\det({}^t(U-\lambda I_n))=\det(U-\lambda I_n)).
\]

Soit $a$ une racine de $\chi_u=\chi_{{}^tu}$ et $\varphi\ne0$, une forme linéaire, vecteur propre de ${}^tu$, on sait que $H=\operatorname{Ker}\varphi$ est un hyperplan de $E$ stable par $u$, (Théorème 10.16).

Si $\bar u$ est induit par $u$ sur $H$, et si $\widetilde U$ est la matrice de $\bar u$ dans une base $C$ de $H$, en prenant un vecteur $\varepsilon\notin H$, $C\cup\{\varepsilon\}$ est une base de $E$ dans laquelle la matrice de $u$ est
\[
U'=
\begin{pmatrix}
\widetilde U&\begin{matrix}\alpha_1\\\alpha_2\\\vdots\\\alpha_{n-1}\end{matrix}\\
0\ \cdots\ 0&\alpha_n
\end{pmatrix}.
\]
Donc
\[
\chi_u(\lambda)=\det(U'-\lambda I_n)
=\det(\widetilde U-\lambda I_{n-1})(\alpha_n-\lambda)
=\chi_{\bar u}(\lambda)(\alpha_n-\lambda).
\]
Comme $\chi_u$ est scindé, il en résulte que $\chi_{\bar u}$ est scindé aussi, donc dans $H$ il existe une base $C'$ dans laquelle la matrice $\widetilde U'$ de $\bar u$ est triangulaire supérieure.

La matrice $U''$ de $u$ dans la base $C'\cup\{\varepsilon\}$ devient donc
\[
U''=
\begin{pmatrix}
\widetilde U'&\begin{matrix}\alpha'_1\\\alpha'_2\\\vdots\\\alpha'_{n-1}\end{matrix}\\
0\ \cdots\ 0&\alpha_n
\end{pmatrix}:
\]
elle est bien triangulaire supérieure puisque la matrice $\widetilde U'$ l'est.

Il est à remarquer que $\alpha_n$ est inchangé, car le spectre de $\widetilde U$ est celui de $\widetilde U'$, et celui de $U'$ est celui de $U''$, en appelant ici spectre les valeurs propres comptées avec leur multiplicité.

On a bien justifié le résultat. Il est à noter que l'emploi du transposé permet de faire une récurrence en passant directement de l'ordre $n-1$ à l'ordre $n$. \bookend

\medskip\noindent\textsc{Corollaire 10.32.} --- \emph{Tout endomorphisme $u$ de $E\simeq\mathbb C^n$ est trigonalisable.}

\medskip\noindent\textsc{Corollaire 10.33.} --- \emph{Toute matrice carrée complexe, est semblable à une matrice triangulaire.}

C'est évident puisque $\mathbb C$ est algébriquement clos.

Vu la difficulté de la démonstration, on peut se demander s'il s'agit là d'une notion intéressante. Et oui! Par exemple on peut justifier que

\medskip\noindent\textsc{Théorème 10.34.} --- \emph{Soit $\{\lambda_1,\ldots,\lambda_n\}$ le spectre de $u$ endomorphisme de $\mathbb C^n$. Si $P$ est un polynôme, le spectre de $P(u)$ est $\{P(\lambda_1),\ldots,P(\lambda_n)\}$.}

On appelle ici spectre de $u$ l'ensemble de ses valeurs propres comptées avec leur multiplicité : si $\lambda_i$ est valeur propre de multiplicité $\alpha_i$ en tant que racine du polynôme caractéristique, elle est répétée $\alpha_i$ fois.

Soit
\[
P(X)=\sum_{k=0}^pa_kX^k
\]
un polynôme de $\mathbb C[X]$. Avec $u^k=u\circ\cdots\circ u$, ($k$ fois), on note $P(u)$ l'endomorphisme
\[
\sum_{k=0}^pa_ku^k,
\]
(avec $u^0=\mathrm{id}_E$).

Le spectre de $u$ étant $\{\lambda_1,\ldots,\lambda_n\}$ et $u$ étant trigonalisable, il existe une base de $\mathbb C^n$ dans laquelle la matrice $U$ de $u$ est triangulaire supérieure avec les $\lambda_i$ sur la diagonale :
\[
U=
\begin{pmatrix}
\lambda_1&u_{12}&u_{13}&\cdots&u_{1n}\\
0&\lambda_2&u_{23}&\cdots&u_{2n}\\
0&0&\lambda_3&\cdots&u_{3n}\\
\vdots&\vdots&\ddots&\ddots&\vdots\\
0&0&0&\cdots&\lambda_n
\end{pmatrix}.
\]
Il est alors facile de justifier, (par récurrence sur $k$) que dans cette base, la matrice $U^k$ sera encore triangulaire, avec sur la diagonale, $\lambda_1^k,\lambda_2^k,\ldots,\lambda_n^k$ d'où finalement $P(u)$ aura pour matrice dans cette base une matrice triangulaire supérieure avec, sur la diagonale, les valeurs $P(\lambda_i)$.

Mais comme le spectre d'une matrice triangulaire est formé des éléments diagonaux (calcul évident du polynôme caractéristique) le spectre de $P(u)$ est bien formé des $P(\lambda_i)$. \bookend

\medskip\noindent\textsc{Remarque 10.35.} --- Soit $u$ un endomorphisme de $\mathbb R^n$, son polynôme caractéristique $\chi_u(X)\in\mathbb R[X]$ et peut avoir des racines imaginaires non réelles, $\lambda_k$. Par contre, il se peut que $P(\lambda_k)$ devienne réel. Dans ce cas, si on note $\{\lambda_1,\ldots,\lambda_n\}$ les zéros de $\chi_u$, comptés avec leur multiplicité, le spectre de $u$, (dans $\mathbb R$) est l'ensemble des $\lambda_k$ réels, et si $P(X)\in\mathbb R[X]$, celui de $P(u)$ est l'ensemble des $P(\lambda_j)$ réels, même si les $\lambda_j$ sont complexes.

Par exemple, dans $\mathbb R^2$ si $u$ est une rotation d'angle $\frac\pi2$, de matrice
\[
\begin{pmatrix}0&-1\\1&0\end{pmatrix},
\]
dans une base orthonormée de $\mathbb R^2$ euclidien, le spectre de cette matrice est $\{i,-i\}$ dans $\mathbb C$, donc vide dans $\mathbb R$, mais le spectre de $u^2$, rotation d'angle $\pi$, donc symétrie, est devenu réel, c'est $\{-1,-1\}$, celui de $u^4$ ($=\mathrm{id}_E$) est $\{1,1\}$.

\medskip\noindent\textsc{Remarque 10.36.} --- Ce qui vient d'être fait pour $\mathbb R$ et $\mathbb C$ est valable pour un corps $K$ quelconque en utilisant $\overline K$, clôture algébrique de $K$ par exemple, (Théorème 7.70) pour obtenir un corps dans lequel le polynôme caractéristique $\chi_u$ d'une matrice de $M_n(K)$ est scindé. On passe dans $\overline K$ pour justifier l'existence d'une forme commode de la matrice, mais si $P\in K[X]$, le calcul de $P(U)$ s'effectue dans l'espace vectoriel $M_n(K)$ des matrices carrées d'ordre $n$ sur $K$.

\medskip\noindent\textsc{Théorème 10.37.} --- \emph{Soient $u$ et $v$ deux endomorphismes trigonalisables qui commutent. Il existe une base commune de $K^n$, de trigonalisation.}

On va employer un résultat, utile en soi.

\medskip\noindent\textsc{Lemme 10.38.} --- \emph{Soient $f$ et $g$ deux endomorphismes de $E$, vectoriel de dimension finie, qui commutent, et qui ont des polynômes caractéristiques scindés. Alors ils ont un vecteur propre non nul, commun, au moins.}

Soit $\lambda$ un zéro de $\chi_f(X)$ polynôme caractéristique de $f$, et $E_\lambda$ le sous-espace propre associé. Comme $f$ et $g$ commutent, $E_\lambda$ est stable par $g$, (Théorème 10.14). Soit $\bar g$ induit par $g$ sur $E_\lambda$, $C$ une base de $E_\lambda$, $D$ une base d'un supplémentaire et $B=C\cup D$ une base de $E$. La matrice de $g$ dans $B$ est du type
\[
V=\begin{pmatrix}M&N\\0&P\end{pmatrix},
\]
matrice bloc, avec $M$ matrice de $\bar g$ dans la base $C$ de $E_\lambda$.

Donc, si $p=\dim E_\lambda$, on a, (calcul par blocs) :
\[
\chi_g(X)=\det(V-XI_n)=\det(M-XI_p)\det(P-XI_{n-p}),
\]
avec $\det(M-XI_p)=\chi_{\bar g}(X)$, polynôme caractéristique de $\bar g$.

Ce polynôme caractéristique de $\bar g$ divise donc celui de $g$. Or ce dernier est scindé, donc $\chi_{\bar g}(X)$ est scindé. Soit $\mu$ un zéro dans $K$ de $\chi_{\bar g}$, et $a$ un vecteur propre, $\ne0$, de $\bar g$ pour $\mu$, on a $g(a)=\mu\cdot a=\bar g(a)$, et comme $a\in E_\lambda$ on a aussi $f(a)=\lambda a$ : le vecteur $a$ est vecteur propre non nul, commun pour $f$ et $g$. \bookend

Justifions le théorème 10.37.

Si $u$ et $v$ commutent, il en est de même de ${}^tu$ et de ${}^tv$, (Théorème 6.101). De plus $u$ et $v$ étant trigonalisables, \BellaCorrection{leurs polynômes caractéristiques}{leur polynômes caractéristiques} sont scindés, or $\chi_u=\chi_{{}^tu}$, (on transpose les matrices). Donc ${}^tu$ et ${}^tv$ sont deux endomorphismes du dual $E^*$ qui commutent ayant des polynômes caractéristiques scindés : d'après le lemme 10.38 ils ont des vecteurs propres communs.

Soit $\varphi$ une forme linéaire non nulle, vecteur propre de ${}^tu$ et ${}^tv$. L'hyperplan $H=\operatorname{Ker}\varphi$ est stable par $u$ et par $v$, (Théorème 10.16). Soient $u'$ et $v'$ induits par $u$ et $v$ sur $H$ : ils commutent, comme $u$ et $v$ en fait, et leurs polynômes caractéristiques sont scindés car ils divisent $\chi_u$ et $\chi_v$ qui le sont, justification analogue à celle du lemme 10.38. On comprend que la justification s'achève par récurrence.

Si $n=1$, toute base de $E=K$ est base commune de trigonalisation.

Si le résultat est vrai pour les espaces de dimension $n-1$, soit $E$ de dimension $n$, $H$ construit comme précédemment, $a$ tel que $E=H\oplus K\cdot a$.

Pour $u'$ et $v'$ induits par $u$ et $v$ sur $H$, l'hypothèse de récurrence s'applique : il existe une base $B_1$ de $H$ dans laquelle les matrices $U'$ et $V'$ de $u'$ et $v'$ sont triangulaires supérieures. Mais alors, dans $B=B_1\cup\{a\}$, les matrices $U$ et $V$ de $u$ et $v$ sont
\[
U=\begin{pmatrix}U'&\begin{matrix}\alpha_1\\\alpha_2\\\vdots\\\alpha_{n-1}\end{matrix}\\0\ \cdots\ 0&\alpha_n\end{pmatrix}
\qquad\text{et}\qquad
V=\begin{pmatrix}V'&\begin{matrix}\beta_1\\\beta_2\\\vdots\\\beta_{n-1}\end{matrix}\\0\ \cdots\ 0&\beta_n\end{pmatrix},
\]
elles sont triangulaires supérieures puisque $U'$ et $V'$ le sont. \bookend

\medskip\noindent\textsc{Corollaire 10.39.} --- \emph{Soient $U$ et $V$ deux matrices carrées d'ordre $n$ sur $K$ qui commutent ayant des polynômes caractéristiques scindés, il existe $P$ régulière telle que $P^{-1}UP$ et $P^{-1}VP$ soient triangulaires.}

Il suffit d'introduire $E=K^n$, une base de $E$, par exemple la base canonique, et de considérer les endomorphismes $u$ et $v$ de matrices $U$ et $V$ dans la base $B$. \bookend
